\section{Introdução}
A modelagem matemática desempenha um papel fundamental nas finanças quantitativas ao fornecer ferramentas rigorosas, como a teoria da probabilidade e o cálculo estocástico, para representar a evolução temporal de variáveis sujeitas à incerteza de mercado. A importância dessa formalização torna-se ainda mais evidente na análise de  derivativos financeiros, cujos valores presentes são determinados por outros ativos financeiros, chamados de ativos subjacentes.
Entre os derivativos mais negociados globalmente, o mercado de opções ocupa uma posição de destaque \cite{fia2026options}. Por concederem ao seu detentor o direito, não a obrigação, de negociar um ativo a um preço fixado em uma data futura, as opções apresentam uma estrutura não linear de ganhos que as torna essenciais para a transferência e mitigação de riscos em carteiras de investimento. Segundo dados divulgados pela Futures Industry Association (FIA, 2026), aproximadamente 28,61 bilhões de contratos de opções foram negociados mundialmente no primeiro trimestre de 2026, em comparação com 9,74 bilhões de contratos futuros. Esses números evidenciam a relevância desse segmento no mercado internacional de derivativos negociados em bolsa. Essa flexibilidade contratual, associada à incerteza em relação a evolução futura dos preços dos ativos subjacentes, torna a precificação de opções um problema relevante para finanças quantitativas. Nesse quadro, o modelo desenvolvido por Black e Scholes \cite{blackscholes1973} estabeleceu uma fórmula fechada para determinar o preço teórico das opções de compra e venda europeias sob determinadas condições de mercado.

O modelo tem como base a representação do preço do ativo como um processo estocástico através do movimento browniano geométrico, assumindo parâmetros como volatilidade e taxa livre de juros constantes. Ao partir do objetivo da construção de uma carteira de proteção, ou também conhecido como carteira de \textit{Hedge}, continuamente ajustada e da ausência de arbitragem, obtém-se uma equação diferencial parcial em que a solução permite determinar o valor teórico do derivativo. Nesse cenário, a precificação passa a ser fundamentada em uma relação matemática entre as características do contrato, como o tempo até o seu fim, denominado de maturação, e toda a dinâmica envolvida na precificação desse ativo.

Para opções europeias de compra e venda, a formulação fornece expressões analíticas que relacionam o preço da opção ao valor corrente do ativo, ao preço de exercício, ao tempo restante até o vencimento, à taxa livre de risco e à volatilidade. Além de fornecer uma expressão para o preço teórico das opções, o modelo de Black-Scholes permite investigar a influência de seus parâmetros sobre o valor do derivativo, por meio das sensibilidades financeiras conhecidas como Gregas. A existência de uma formulação analítica possibilita, ainda, examinar as propriedades matemáticas da função de precificação, tornando o modelo uma referência para a avaliação de métodos computacionais de aproximação.

A existência dessa solução analítica permite utilizar o modelo como referência para a investigação de métodos computacionais capazes de aproximar relações matemáticas entre variáveis do problema. Nesse quadro, pode-se destacar a regressão simbólica como uma técnica avançada em aprendizado de máquina que utiliza programação genética e algoritmos evolutivos para identificar expressões matemáticas explícitas a partir dos dados, além de buscar aproximações numéricas que possibilitam investigar estruturas matemáticas e relações subjacentes aos dados. A regressão simbólica foi projetada para inferir expressões matemáticas interpretáveis orientada a dados, diferente de métodos cuja estrutura é previamente estabelecida. Entre as ferramentas disponíveis para essa finalidade, destaca-se o algoritmo PySR, que emprega técnicas da já mencionada busca evolutiva e aprendizado por árvores de decisão para identificar expressões simbólicas capazes de representar as relações presentes nos dados. Sua aplicação a uma função analítica conhecida, como a de Black-Scholes, permite investigar não apenas a capacidade de aproximação dos preços, mas também a correspondência entre as expressões identificadas e as propriedades matemáticas do modelo original.

\subsection{Tema}

\subsection{Justificativa}

\subsection{Delimitação}

\subsection{Objetivos}
Investigar a capacidade da regressão simbólica, por meio do algoritmo PySR, de aproximar a função de precificação de Black-Scholes, avaliando a precisão preditiva, a complexidade simbólica e a recuperação de propriedades matemáticas do modelo.

\subsection{Descrição resumida da metodologia}

\subsection{Descrição}
